\section{Introduction}
A service provider can improve its operating policy and still offer an
unacceptable contract. A revision of service tiers changes customer payments,
future service obligations, and switching expenditures. When a customer may
leave at a review, an improvement at the root of the planning tree does not
suffice: every promised continuation must remain acceptable. The relevant
operations research question is which expansions of an already optimized,
accepted contract create value, and how that value can be implemented and
verified.

We develop an accepted-service expansion theory on finite history trees.
Physical decisions and contractual tiers are publicly observable; tariffs and
the demand law are fixed before optimization. The customer accepts a contingent
protocol subject to physical and continuation commitments. There is no private
type or hidden effort. The restricted comparator already optimizes its
permitted decisions under those same commitments. Adaptation therefore does
not earn a spurious advantage by being compared with an arbitrary tier or a
state-blind inventory rule.

For the positive-payment tree specialization, every accepted expansion around a
restricted contract has unique continuation-transfer coordinates, whose lower
bounds are the contract's actual participation slack. Cumulative continuation
prices and switching capacities identify whether a feasible improvement exists.
In the cap-matching case, a critical-friction characterization and a two-pass
tension repair produce a profitable accepted deviation or a verified upper
threshold. For a general optimized comparator, the translated coordinates retain
nonconstant tiers, unused acceptance slack, and boundary decisions. A quadratic
objective gives an exact benchmark-relative expansion formula. The distinction
is substantive: forcing nonnegative transfers around a slack comparator can
exclude the economically optimal accepted revision.

The second contribution connects these accepted-control objects to the value
of the implemented decision. A primal--dual component identity separates
resource-price mismatch, switching tension, box optimality, and participation
slack. It shows what price-only repair can and cannot accomplish; complete dual
repair attains the true regret floor without a strict-feasibility assumption.
A feasible gain witness and a certified restricted upper bound then establish
economic value against the optimized contract, rather than merely against an
announced outside option. These are service-control and certification
consequences of stated convex tools, not claims to replace general duality,
network-flow theory, or dynamic programming.

Neural Differential Utility provides one implementation route within this
framework. Under explicit strong convexity and smooth resource coupling, a
partial value gradient in specified reward coordinates supplies a useful
resource statistic. Direct-price regression and structure-aware classical
continuation use the same accepted geometry and remain full comparators. We
retain the global bridge, inexact-response bounds, scalar and verified-face
developments, and all original experiments. The geometric source of expansion
value, the value recovered by a proposal, and the cost of recovering it are
separate scientific questions. No theorem assumes that a neural predictor is
necessary, and no unfavorable comparison is removed.

The experiments measure accepted gains against reoptimized time-only
amendments and computation against lifted primal--dual continuation. Repeated
query measurements charge optimizer-generated labels and their exact audits, fitting, prediction, auditing,
polishing, and fallback. A designed joint-coefficient uncertainty study further
tests acceptance and economic certificates when the objective and commitments
change together. Direct interior-model optimization quantifies the conservatism
of its common outer comparator. Truth-in-set certification, sensitivity to a
designed perturbation family, statistical calibration, and field robustness are
distinct; the last two are not inferred from synthetic stress tests.

